%=====================================================================
\clearpage
\section{Erweiterung ab Version 2.22: CAD-Import für 1D-Gitter (M23)}\label{sec:cad1d}
%=====================================================================
Der CAD-Slicer (M14, Abschnitt~\ref{sec:cad}) schneidet Körper in der $(x,z)$-Zelle und ist deshalb auf 2D-Gitter
beschränkt (v2.21). Liniengitter sind aber der häufigste Fall: ein in $z$ extrudiertes Profil, z.\,B. ein Trapez mit
exakten Kreisbögen an Kopf und Fuß aus \verb|nurbforge.fillet_corners|. Als 2D-Gitter lässt sich so ein Körper nicht
darstellen: Über eine $z$-Periode extrudiert, berührt er sein Bild, und der Schnitt ist ein Streifen um die Zelle,
keine geschlossene Kontur (\verb|PERIODIC_OVERLAP|, v2.21). \verb|profile1d| deckt nur eine einzige Grenzfläche
$y=g(x)$ ab, ohne Überhänge und ohne mehrere Materialien. M23 ergänzt den Slicer um einen 1D-Pfad: $z$-invariante
Körper werden \emph{einmal} in einer Ebene $z=z_c$ geschnitten. Das liefert das Profil in der $(x,y)$-Ebene. Daraus
entstehen dieselben Bänder wie in M14, aber mit 1D-Schichten.

\subsection{Annahmen und Festlegungen}
\begin{enumerate}[nosep]
  \item Gleiche \verb|cad|-Schicht, gleiche Schlüssel und gleiche Prüfcodes wie M14. Der 1D-Pfad wird allein durch
        \verb|"dimension": "1D"| gewählt. Der 2D-Pfad bleibt unverändert und bitgleich (alle M14-Tests, Golden Files).
  \item Jeder Körper muss $z$-invariant sein, sonst Befund \verb|NOT_Z_INVARIANT| (neu). Die Extrusionslänge in $z$ ist
        beliebig und geht in keine Prüfung ein, auch nicht in \verb|UNIT_SUSPECT|.
  \item Nurbforge wird nur zum Lesen, für die Prüfungen, die kritischen Höhen und den \emph{einen} Schnitt bei $z_c$
        gebraucht. Alles danach (Linienschnitt, Intervalle, Normalen, $Q(y)$) ist FG-Geometrie auf
        \verb|geometry::BezierSegment| und hängt nicht von Nurbforge ab. So lässt es sich ohne
        \verb|FG_ENABLE_NURBFORGE| testen. Nach dem Stand von v0.6.0 ist keine Vorleistung in Nurbforge nötig
        (\verb|section| mit beliebiger Ebene, \verb|critical_heights|, \verb|is_extrusion_band| mit Richtung $z$).
  \item Neuer interner Schichttyp \verb|Sectioned1DLayer|, wie \verb|SectionedLayer| nicht in JSON verfügbar. Ein
        eigener JSON-Typ für Profile aus Bézierkurven wäre möglich, ist aber nicht Teil von M23 (keine Funktionen
        auf Vorrat).
  \item Normalenfeld ohne Vorzeichen (Projektorform): erlaubt sind dieselben Faktorisierungen wie bei
        \verb|Contour2D|. Standard ist \verb|FFF_sym| (isotrop) bzw.\ \verb|FFF_B| (anisotrop).
\end{enumerate}

\subsection{Meilenstein und Abhängigkeiten}
\begin{center}
\small
\begin{tabular}{p{1.0cm}p{6.6cm}p{3.0cm}p{3.8cm}}
\toprule
M & Inhalt & setzt voraus & prüfbar in der Sandbox \\
\midrule
M23 & CAD-Import für 1D-Gitter: $z$-Invarianzprüfung, Profilschnitt bei $z_c$, Linienschnitt $\to$ Intervalle mit
      Normalen, \verb|Sectioned1DLayer| (RK, RCWA, Tschebyschow, ASR, Nahfeld, Gradienten), Bericht und Vorschau
      & M14, M17, M19 & ja (Nurbforge per \verb|FetchContent|) \\
\bottomrule
\end{tabular}
\end{center}

\subsection{Datenfluss}
{\footnotesize\begin{verbatim}
STEP --read_step_named--> Körper --Einheit/Achsen--> FG-Koordinaten (µm)            (wie M14)
  --is_extrusion_band(z)--> NOT_Z_INVARIANT?                                           (neu)
  --section(z = z_c)--> Profilschleifen je Körper in (x, y), Bézierstücke              (neu, einmal)
  --critical_heights(y)--> Bänder [y_k, y_k+1]                                         (wie M14)
  --je Band: homogen | lamellar1d (Extrusion in y, ein Linienschnitt) | Sectioned1DLayer
  --Linienschnitt y = const--> Intervalle (Material) + Normalen an den Bruchstellen    (neu, FG)
\end{verbatim}}

\subsection{$z$-Invarianz und Schnittebene}
\begin{itemize}[nosep]
  \item Prüfung je Körper mit dem getrimmten Hüllquader $[z_0,z_1]$ (v2.21):
        \verb|nf::is_extrusion_band(shell, {0,0,1}, z0, z1)| muss wahr sein. Die Deckel liegen auf den Bandgrenzen und
        zählen deshalb nicht. Ist die Prüfung falsch, gibt es den Befund \verb|NOT_Z_INVARIANT| mit der Körpermitte als Ort.
  \item Schnittebene: \verb|"z_cut_um"| in FG-Koordinaten. Standard ist die Mitte des gemeinsamen $z$-Bereichs aller
        Körper. Ist dieser Bereich leer oder liegt $z_c$ nicht im Inneren jedes Körpers, folgt \verb|OUT_OF_RANGE|.
        Wegen der $z$-Invarianz hängt das Ergebnis nicht von $z_c$ ab (Test V24e).
  \item Schnitt: \verb|nf::section(shell, {{0,0,z_c},{0,0,1}})|, \verb|s_axis| $=\hat{\mathbf x}$. Die Kurven werden wie in
        M14 in Bézierstücke zerlegt, aber auf $(x,y)$ abgebildet. Die Orientierung der Schleifen wird nicht
        verwendet, das Innere folgt aus der Paritätsregel (Abschnitt~\ref{sec:cad1d_line}). Löcher (Hohlräume) brauchen
        deshalb keine eigene Behandlung.
  \item \verb|cell_origin_um| ist in 1D eine Zahl $x_0$ oder \verb|[x0]|. Ein zweiter Eintrag ist ein Fehler. Standard
        ist das Minimum der getrimmten Hüllquader in $x$.
\end{itemize}

\subsection{Linienschnitt: Intervalle, Materialien, Periodizität}\label{sec:cad1d_line}
Für eine Höhe $y$ im Inneren eines Bandes werden je Körper $b$ die Schnittpunkte der Geraden $y=\text{const}$ mit allen
Bézierstücken seines Profils bestimmt. Methode: Bernstein-Unterteilung, bis die konvexe Hülle die Gerade nur noch
einmal kreuzt, dann Newton auf $y(t)-y$ für rationale Stücke; Toleranz $10^{-14}$ der Zellgröße. An jedem
Schnittpunkt $x$ wird die Profilnormale
\begin{equation}
  \mathbf n=\frac{(\dot y(t),\,-\dot x(t),\,0)}{\sqrt{\dot x^2+\dot y^2}}\label{eq:cad1d_n}
\end{equation}
festgehalten (Vorzeichen egal). Bei $z$-invarianten Körpern ist sie die 3D-Flächennormale. Da im Inneren eines Bandes
nie auf einer kritischen Höhe geschnitten wird, gibt es keine tangentialen Schnitte. Die Zahl der Schnittpunkte ist
daher konstant und gerade.
\begin{enumerate}[nosep]
  \item Inneres je Körper nach der Paritätsregel: Schnittpunkte nach $x$ sortieren (ohne Rückfaltung). Innen liegt
        $[x_1,x_2]$, $[x_3,x_4]$, \dots
  \item Rückfaltung in die Zelle $[x_0,x_0+d_x)$, danach relativ zu $x_0$. Überlappen sich zwei Intervalle
        \emph{desselben} Körpers nach der Rückfaltung um mehr als die Toleranz, ist das \verb|PERIODIC_OVERLAP|. Mit
        \verb|"periodic_union": true| werden sie vereinigt (in 1D eine Intervall-Vereinigung, keine Boolesche Operation in 3D).
  \item Berühren sich Intervalle desselben Körpers nur (gemeinsamer Endpunkt innerhalb der Toleranz, z.\,B. eine Schicht
        über die ganze Periode oder ein Körper, der über den Zellrand ragt), werden sie ohne Befund verschmolzen. Die
        Naht ist keine Materialgrenze, ihre Normale entfällt. Das ist der wesentliche Unterschied zu 2D (v2.21).
  \item Überlappen sich Intervalle \emph{verschiedener} Körper: \verb|BODY_OVERLAP|. Bei \verb|"overlap": "priority"|
        wird die Intervall-Differenz $B_i\leftarrow B_i\setminus\bigcup_{j<i}B_j$ je Schnitt gebildet; eine Boolesche
        Operation in 3D entfällt. Berühren sie sich nur, bleibt die Bruchstelle (echte Materialgrenze).
  \item Ergebnis: Bruchstellen $0=\xi_0<\xi_1<\dots<\xi_K=d_x$ und Materialien je Intervall. Der Hintergrund füllt die
        Lücken, Intervalle der Länge $<$ Toleranz entfallen. Liegen zwei benachbarte Intervalle im selben Material,
        entfällt die Bruchstelle. Jede verbleibende Bruchstelle $\xi_k$, $0<k<K$ (und eine Bruchstelle bei $0\equiv d_x$,
        wenn dort das Material wechselt) trägt ihre Normale $\mathbf n_k$ aus \eqref{eq:cad1d_n}.
\end{enumerate}
Die Überlappungsprüfung des Slicers verwendet genau diese Funktion an drei Höhen je Band ($0{,}2$, $0{,}5$, $0{,}8$), wie in
M14. Prüfung und Löser können damit nicht mehr auseinanderlaufen.

\subsection{Normalenfeld im Linienschnitt}\label{sec:cad1d_nf}
Seien $\xi_1<\dots<\xi_M$ die Materialgrenzen eines Schnitts (zyklisch, $\xi_{M+1}=\xi_1+d_x$) mit Winkeln
$\theta_k=\operatorname{atan2}(n_{k,y},n_{k,x})\bmod\pi$. Auf $[\xi_k,\xi_{k+1}]$ gilt
\begin{equation}
  \theta(x)=\theta_k+\frac{x-\xi_k}{\xi_{k+1}-\xi_k}\,\delta_k,\qquad
  \delta_k=\theta_{k+1}-\theta_k\ \text{reduziert auf}\ (-\tfrac\pi2,\tfrac\pi2],\qquad
  \mathbf N(x)=(\cos\theta,\,\sin\theta,\,0).\label{eq:cad1d_N}
\end{equation}
$\mathbf N\mathbf N^{\mathsf T}$ ist stetig und an jeder Grenze exakt die Grenzflächennormale. Das ist die Bedingung der
FFF (Abschnitt~\ref{sec:fff}). Wegen $\cos^2\theta=\tfrac12(1+\cos2\theta)$ und $\cos\theta\sin\theta=\tfrac12\sin2\theta$
mit linearem $\theta$ haben die Koeffizienten von $N_iN_j$ geschlossene Form: Je Intervall ist es das Integral von
$e^{\pm2i\theta(x)}e^{-iK_px}$, also eine Exponentialfunktion mit linearer Phase. Umsetzung über
\verb|fourier::piecewiseSmoothCoefficients| mit Gauß-Legendre (Test gegen die geschlossene Form auf $10^{-14}$).
Sonderfälle: $M=0$ (homogener Schnitt) braucht kein Feld. Sind alle $\theta_k=0$ (senkrechte Wände), ist
$\mathbf N=\hat{\mathbf x}$, und \verb|FFF_sym| fällt mit der Li-Regel für Lamellen zusammen (Test V24a).

Für ein einwertiges Profil $y=g(x)$ liefert \verb|profile1d| ein anderes Feld \eqref{eq:N1d}: $\mathbf N=\hat{\mathbf y}$
auf Kopf und Grund statt interpoliert zwischen den Wänden. Beide Felder sind an den Grenzen gleich. Die Koeffizienten
von $\varepsilon$ und $1/\varepsilon$ stimmen exakt überein, $Q$ nur bei Laurent (V24b). Die Effizienzen
konvergieren gegen denselben Grenzwert.

\subsection{Schichttyp \texttt{Sectioned1DLayer}}
{\footnotesize\begin{verbatim}
struct LineSection { std::vector<double> breaks;            // 0 = xi_0 < ... < xi_K = dx
                     std::vector<std::string> materials;    // K Einträge
                     std::vector<Vec3D> normals; };         // je Materialgrenze (Projektorform)
class LineSectionProvider {                                 // M23: Profil aus CAD; Tests: analytisch
 public: virtual LineSection section(double y_um) const = 0; // y relativ zur Schichtunterkante
         virtual double thickness_um() const = 0; };
struct Sectioned1DLayer { std::shared_ptr<const LineSectionProvider> provider;
                          std::string background; int slices = 0; int chebyshevNodes = 0; };
\end{verbatim}}
\begin{itemize}[nosep]
  \item Modell \verb|detail::Sectioned1DModel|: Es verlangt einen 1D-Stapel (ein 2D-Stapel ist ein Fehler, wie
        umgekehrt bei \verb|SectionedLayer|). $Q(y)$ wird aus \verb|isoGrids| bzw.\ \verb|anisotropicQ| über den
        Intervallen und dem Feld \eqref{eq:cad1d_N} gebildet, genau wie bei \verb|ProfileModel|, nur mit Bruchstellen und
        Normalen vom Provider. Cache je $y$ (bitgenau) und thread-sicher wie bei \verb|SectionedModel|.
  \item Tschebyschow-Interpolation in $y$ wie in M14. Der Interpolant (Knoten, Gewichte, Fehlerkontrolle $K$ gegen $2K$)
        wird aus \verb|SectionedModel| in eine gemeinsame Hilfsklasse ausgelagert, nicht kopiert.
  \item RCWA: Scheibenmitten, die Treppe besteht aus Lamellen. RK: Stufenknoten. Der Kontaminationsschutz
        \eqref{eq:slicecrit} gilt unverändert.
  \item M17 (ASR in $x$): Bruchstellen über \verb|mapIntervals| abbilden, die Normale an der physikalischen Stelle
        $x=f(X)$ auswerten, wie bei \verb|ProfileModel|.
  \item Nahfeld: \verb|epsAt| aus dem Schnitt bei $y$, \verb|normalAt|/\verb|NN| aus \eqref{eq:cad1d_N}.
  \item M19: \verb|adjointSlice| wie bei \verb|ProfileModel| (Intervalle, Indikatoren, Normalenfeld, anisotrop über
        \verb|anisoAdjointSlice|), nicht bei Tschebyschow-Interpolation. Formparameter laufen über
        \verb|StackParameter| (Neuaufbau aus der STEP-Datei), wie im Test \verb|CAD_M19|. Ableitung (v2.23): zentrale
        Differenz der Koeffizientengitter ($\varepsilon$, $1/\varepsilon$, $N_iN_j$, $N_i$ bzw.\ die anisotropen
        Funktionen) der Schnitte derselben Höhe bei $\theta\pm\delta$ (\verb|AdjointSlice::gridsOfSection|); das
        Normalenfeld \eqref{eq:cad1d_N} wird also mitdifferenziert.
  \item PML (M18): nicht mit \verb|Sectioned1DLayer| (klare Fehlermeldung).
\end{itemize}

\subsection{Bandklassifikation in 1D}
Die Bänder entstehen wie in M14 aus \verb|nf::critical_heights| in $y$. Die Klassifikation stützt sich auf den
Linienschnitt in Bandmitte:
\begin{itemize}[nosep]
  \item ein Material über die ganze Periode $\to$ \verb|homogeneous|;
  \item alle Körper im Band erfüllen \verb|is_extrusion_band| in $y$ $\to$ \verb|lamellar1d| aus diesem Schnitt (exakt);
  \item sonst \verb|Sectioned1DLayer| mit einem \verb|ProfileLineProvider| (Profilschleifen des Bandes,
        Hintergrund, Toleranz).
\end{itemize}
\verb|FEATURE_SMALL| warnt in 1D, wenn ein Intervall in Bandmitte schmaler als zwei Rasterabstände
$2d_x/\max(\texttt{grid},64)$ ist wie in 2D (v2.23; $N_x$ ist dem Slicer nicht bekannt). Im Bericht heißen die
Bandarten \verb|homogeneous|, \verb|lamellar1d| und \verb|sectioned1d|. Die Flächenanteile werden zu Längenanteilen
(bezogen auf $d_x$).

\subsection{Konfiguration, Bericht, Vorschau}
{\footnotesize\begin{verbatim}
{"dimension": "1D", "period_um": {"x": 0.3}, "truncation": {"Nx": 30},
 "layers": [{"type": "cad", "file": "line_grating.step", "unit": "um", "axes": {"y": "+Z", "x": "+X"},
             "z_cut_um": 0.15, "cell_origin_um": 0.0, "background": "superstrate",
             "bodies": [{"name": "Line", "material": "si"}],
             "slicing": {"rcwa_slices": 40, "y_interpolation": "chebyshev"}}]}
\end{verbatim}}
\begin{itemize}[nosep]
  \item \verb|fg.cad.check| / \verb|--cad-report|: zusätzlich \verb|"z_cut_um"| und \verb|"dimension": "1D"| im Bericht.
  \item \verb|fg.cad.preview| in 1D: SVG des Profils in der $(x,y)$-Ebene über $[0,d_x]$ mit den Bildern bei $\pm d_x$,
        Bandgrenzen als waagrechte Linien. Das Argument \verb|y_um| entfällt.
  \item \verb|llm_tool_spec()| nennt den 1D-Fall und \verb|NOT_Z_INVARIANT|.
\end{itemize}

\subsection{Tests und Abnahme (M23)}
Der Provider-Teil wird ohne Nurbforge geprüft (\verb|tests/unit/test_line_section.cpp|, analytische Provider). Der
CAD-Teil steht in \verb|tests/cad/test_cad.cpp| (Label \verb|validation|), STEP-Dateien werden zur Laufzeit geschrieben.
\begin{itemize}[nosep]
  \item \textbf{V24a} Rechtecklinie (senkrechte Wände): ein Band \verb|lamellar1d|. $Q$ gleich \verb|lamellar1d|
        ($10^{-13}$), Effizienzen RK und RCWA gleich ($10^{-12}$); \verb|FFF_sym| des Providers ist gleich Li ($10^{-13}$).
  \item \textbf{V24b} Trapez (Wandwinkel $15^\circ$): ein Band \verb|sectioned1d|. Mit Laurent stimmen $Q(y)$ an fünf
        Höhen mit \verb|profile1d| \verb|points| überein ($10^{-13}$), ebenso die RK-Effizienzen ($10^{-10}$). Mit
        \verb|FFF_sym| konvergieren beide gegen denselben Grenzwert (Differenz fällt mit $N_x$).
  \item \textbf{V24c} Gerundetes Trapez wie in der Meldung zu v2.21 ($CD=100$\,nm auf halber Höhe, $h=100$\,nm,
        Eckwinkel $105^\circ$, Radius 10\,nm, Kopf konvex, Fuß konkav, exakte Kreisbögen): gegen \verb|profile1d|
        \verb|Function| mit analytischem $g(x)$ (Laurent, RK $10^{-10}$): $10^{-8}$. Mittlerer Querschnitt
        $\frac1h\int\!\text{Länge}(y)\,dy=CD$ auf $10^{-12}$. Normalen \eqref{eq:cad1d_n} auf den Bögen gegen den
        Radiusvektor ($10^{-12}$).
  \item \textbf{V24d} Nicht als \verb|profile1d| darstellbar: hinterschnittenes Trapez mit Deckel aus zweitem Material
        und einem Hohlraum. RK $\approx$ RCWA (Treppe konvergiert), $R+T=1$ ($10^{-10}$, verlustfrei), Reziprozität
        ($10^{-10}$).
  \item \textbf{V24e} Unabhängigkeit: drei Werte von $z_c$ und Extrusionslängen $10^{-1}$ bis $10^{3}$\,\textmu m liefern
        bitgleiche Schichten. Die Verschiebung von \verb|cell_origin_um| ändert nur die Phasen, die Effizienzen auf $10^{-12}$.
  \item \textbf{V24f} Periodizität: Ein Körper „Restschicht + Linie“ ergibt ein homogenes Band unter einem
        \verb|sectioned1d|-Band, gleich zwei getrennten Körpern ($10^{-12}$). Eine Linie über den Zellrand ist gleich der
        verschobenen. Selbstüberlappung gibt \verb|PERIODIC_OVERLAP|, mit \verb|periodic_union| ist das Ergebnis gleich
        dem vereinigten Profil.
  \item \textbf{V24g} Befunde: Zylinder mit Achse $y$ $\to$ \verb|NOT_Z_INVARIANT|. Eine Linie von 1\,mm Länge erzeugt
        kein \verb|UNIT_SUSPECT|. $z_c$ außerhalb $\to$ \verb|OUT_OF_RANGE|. Zwei überlappende Körper $\to$
        \verb|BODY_OVERLAP|, mit \verb|priority| aufgelöst. Ein 2D-Stapel liefert unverändert alle M14-Ergebnisse
        (Golden File \verb|cad_pillar_array| bitgleich).
  \item \textbf{V24h} Gradienten: Wandwinkel und Bogenradius über \verb|StackParameter| (RCWA und RK) gegen zentrale
        Differenzen ($10^{-5}$ relativ). Materialgradient über \verb|adjointSlice| gegen Differenzen ($10^{-6}$).
  \item \textbf{V24i} Python: \verb|fg.cad.check| und \verb|fg.cad.preview| in 1D. Das Beispiel
        \verb|examples/cad/line_grating_1d.json| (STEP aus \verb|scripts/make_cad_examples.py|) bekommt ein neues Golden
        File; bestehende Golden Files ändern sich nicht.
\end{itemize}
Zusätzlich Unit-Tests: Linienschnitt rationaler Bézierstücke (Kreisbogen, Viertelkreis bis an die Tangente) gegen die
geschlossene Form; Paritätsregel mit Loch; Verschmelzen an der Naht; Koeffizienten von \eqref{eq:cad1d_N} gegen die
geschlossene Form.

\subsection{Umsetzung (v2.23)}\label{sec:cad1d_impl}
\begin{itemize}[nosep]
  \item \verb|src/geometry/line_section.cpp|: \verb|lineCrossings| (Bernstein-Koeffizienten $w_j(y_j-c)$,
        Zeichenwechsel, Halbierung, abgesicherter Newton; Nullstellen bei $t=1$ gehören zum nächsten Stück),
        \verb|lineSection| (Paritätsregel, Rückfaltung, Nähte, Überlappungen, \verb|LineSectionFindings| für die
        Prüfung), \verb|normalField(LineSection)| (bei parallelen Grenzen exakt konstant, sonst Gauß-Legendre;
        Test gegen die geschlossene Form $10^{-14}$), \verb|ProfileLineProvider|.
  \item \verb|Sectioned1DModel| / \verb|LineSliceModel| in \verb|layer_model.cpp|; \verb|ChebyshevQ| gemeinsam mit
        \verb|SectionedModel| (Golden Files bitgleich). \verb|Sectioned1DLayer::materials| nennt alle möglichen
        Materialien (TE/TM-Aufspaltung, Wiederverwendung im $\lambda$-Scan).
  \item Bandgrenzen sind kritische Höhen aus \verb|nf::critical_heights|; ein Fußbogen, der $y=0$ tangential berührt,
        liefert die Unterkante als Newton-Extremum (z.\,B.\ $-1{,}5\cdot10^{-12}$\,\textmu m statt $0$). Das ist physikalisch
        bedeutungslos; Tests beziehen Höhen auf die Bandgrenzen des Berichts.
  \item Ergebnisse: V24b RK $<10^{-10}$, RCWA $<10^{-12}$; V24c gerundete Linie gegen \verb|profile1d| \verb|Function|
        $2{,}6\cdot10^{-13}$, Schnittpunkte und Normalen exakt bis auf die STEP-Rundung; V24d RCWA (100 Scheiben) gegen RK
        $2{,}4\cdot10^{-10}$, $R+T=1$ auf $10^{-14}$, Reziprozität $3{,}5\cdot10^{-7}$ bei $N_x=24$ (fallend); V24h
        Gradienten $\le 2\cdot10^{-8}$ relativ.
\end{itemize}

\subsection{Reihenfolge der Umsetzung}
\begin{enumerate}[nosep]
  \item \verb|LineSection|, Linienschnitt und Normalenfeld \eqref{eq:cad1d_N} in \verb|src/geometry/| mit Unit-Tests.
  \item \verb|Sectioned1DLayer| und \verb|Sectioned1DModel| mit analytischem Provider. Tschebyschow-Hilfsklasse aus
        \verb|SectionedModel| auslagern (Golden Files bitgleich).
  \item Slicer: 1D-Zweig in \verb|loadModel|/\verb|checkOverlaps|/\verb|layersOf|, \verb|NOT_Z_INVARIANT|, Konfiguration.
        Die Ablehnung von 1D-Stapeln (v2.21) entfällt.
  \item Bericht, Vorschau, Python, Beispiel, Handbuch Kapitel~12, \verb|docs/cad.md|; V24a--i.
\end{enumerate}
